% ---------------------------------------------------------------------------
% Appendix: error propagation through the chain, and numerical verifications
% (companion to Section sec:chain). All numbers at n = 200; see the scope
% subsection of the main section.
% ---------------------------------------------------------------------------
\section{The chain under truncation: error propagation and verifications}
\label{app:chain}

\subsection{How a truncation at one stage reaches time $T$}
\label{ssec:chain-error}

A truncated chain replaces the exact stage map by its order-$k$ expansion about
the stage's initial state. Three exact structures decide how the local error
of stage $s$ appears at $T$.

\begin{proposition}[The propagator is algebraic]
\label{prop:propagator}
A perturbation of the state at time $s$ evolves by $\dot{\delta\theta}=-H(t)\,\delta\theta$,
$H=J^{\!\top}J/n+\frac1n\sum_iE_i\nabla^2N_i$, so the propagator
$\Phi(T,s)=\mathrm{T}\!\exp(-\!\int_s^TH)$ obeys
$\|\Phi(T,s)\|\le\exp\!\big(\int_s^T\lambda_{\max}(-H)_+\big)\le\exp\!\big(\int_s^T\tfrac1n\|E\|H_2\big)$:
only the residual-curvature part can expand, and its integrand decays with the
residual. [n=200: $\|\Phi\|\le21$ over a full run in the worst direction, i.e.\ an
integrated negative curvature $\ln 21\approx3$, an average of $10^{-3}$ per unit
time; the $H_2$ form of the bound is valid but loose by $e^{13}$, since it
replaces the negative curvature by its worst case $\|E\|H_2/n\approx10^{-2}$.]
\end{proposition}

\begin{proposition}[Three channels]
\label{prop:channels}
A feature error $\delta\Psi_s$ at stage $s$ reaches $T$ as (i) an additive,
undamped feature offset $\Phi(T,s)\delta\Psi_s$ (since $\Psi_T=\Psi_0+\sum_r\delta\Psi_r$),
hence a kernel error linear in the offset while it is smaller than the
features; (ii) a residual error $\delta E(T)=-\int_s^T\mathrm{T}\!\exp(-\!\int_u^TK)\,\delta K(u)E(u)\,du$,
contracted by $K$ in every lifted direction and sourced only by the kernel
error acting on the residual; (iii) a feedback on later stages,
$\delta(\delta\Psi_r)\sim J'[J^{\!\top}\delta E]$ for $r>s$, bounded by
Proposition~\ref{prop:propagator} and of the order of (ii) times the drive.
\end{proposition}

\begin{theorem}[Global error of the order-$k$ chain]
\label{thm:global-error}
If each stage's order-$(k{+}1)$ remainder is at most $\varepsilon$ times its
first-order increment, then
$\|\Psi_T^{\mathrm{chain}}-\Psi_T\|\le\varepsilon\,\Phi_{\max}\sum_s\|\delta\Psi_s\|$ and
$\|E_T^{\mathrm{chain}}-E_T\|\le\frac{2\varepsilon\Phi_{\max}}n\|\Psi_T\|\sum_s\|\delta\Psi_s\|\int_0^T\|E\|+O(\varepsilon^2)$.
The stage errors are proportional to the stage \emph{increments} and add, so
their sum is $\varepsilon$ times the total feature change along the path, not
$\varepsilon$ times the number of stages: the chain is consistent of order $k$
and stable with an algebraic propagator, hence convergent, uniformly in the
number of stages. Fewer, longer stages at higher order and many short stages
at first order both converge; they differ in cost, with
$S_k\sim\varepsilon^{-1/k}\int(\|\Psi^{(k+1)}\|/\|\Psi'\|)^{1/k}$ stages needed.
\end{theorem}
\begin{proof}
Triangle inequality over stages with Proposition~\ref{prop:channels}(i),
Proposition~\ref{prop:propagator} for the propagator, variation of constants
for (ii). The numerical convergence to three digits across a tenfold range of
stage budgets at $k=1$ is this theorem observed.
\end{proof}

\begin{remark}
For proofs the $k=1$, many-stage end suffices: it is the flow itself, and
nothing in this section needs more than first order per stage plus the
propagator bound. For computation the chain is not cheap: long stages need
high-order directional derivatives of the Jacobian at every expansion point,
short stages need of order $10^4$ of them; either way its cost is that of
integrating the dynamics. Its value is the error accounting, which makes
every truncated or reduced version of it a controlled approximation with a
named error term. The constants are finite but not small in exactly two
circumstances: early, where the residual-curvature part of $H$ makes the
propagator grow (the sharpness-reduction phase), and for errors landing in a
slow mode, where (ii) gives no damping and the offset of (i) is later
amplified by the lift of that mode. That is the chain's own statement of
Propositions~\ref{prop:visibility}--\ref{prop:amplification}: what is wrong,
or unknown, in the slow modes early survives and is amplified; what is wrong
in the fast modes is forgotten.
\end{remark}


\subsection{Numerical verifications of the identities and of the leading orders}
\label{app:chain-verif}

\begin{remark}
A numerical realisation of the chain (one eigendecomposition of the
$N\times N$ kernel per stage, the maps of
Proposition~\ref{prop:chain-identities}(c,d) applied to the network's own
units) reproduces the real gradient-flow run to four digits in $\|E\|$ and to
$0.1\%$ in every kernel eigenvalue at $t=8$, tracks it to $1$--$3\%$ at $T$ on
every design tested including depth~8 and width~512, and is converged to
three digits across a tenfold range of stage budgets [n=200]. It is gradient
flow on the sample integrated in large stages; its only use in this section is
as the exact reference against which every reduced statement is checked.
\end{remark}

\begin{remark}[What was measured, n=200]
(a) is verified to within $10\%$ at initialisation for tanh, ReLU, GELU and
softplus; the bulk is $2$--$5\%$ of the core for tanh at unit gain and
$30$--$50\%$ for ReLU (whose local derivative is a $0/1$ pattern), yet $98$--$100\%$
of the first-order change lies in three singular directions for every
activation. Along the path the independence of (a) fails within a few steps,
as it must once the units adapt to the residual, while the rank persists: at
$t=0.3$ the per-unit-mean split gives a bulk ratio of $0.9$ against the
dispersion prediction of $0.05$, and the cumulative change is still $98\%$ in
three directions; the correlation-free bound (b) holds with $r=1$--$2$, i.e.\
the update lies $93\%$ in $3$ and $99.9\%$ in $9$ directions. Cores of
successive stages add coherently while the drive keeps its direction, which
is the cross-layer alignment measured earlier (randomising one layer's
rotation destroys the lift); the bulk-plus-spike spectrum of trained weights
is the random initialisation plus this coherent low-rank sum.
\end{remark}

\begin{remark}[Verification, n=200]
With $\epsilon=1/s$ and the coefficients through second order (including the
second-order weight change and the feedback of the lifted kernel on the lazy
residual), the predicted lift of init eigen-directions agrees with the
measured one to $0$--$16\%$ at output scales $s=2,4,8,16$ for as long as the
measured lift is below order one, including the early \emph{decrease} along
deep-tail directions and its sign change; the truncated series fails sharply
once the lift exceeds one, where Theorem~\ref{thm:radius} says it must. The
leading coefficient $a_1$ changes sign along $\tau$, so no universal power of
$s$ describes the delay to a given lift; the function is predicted, the
exponent does not exist.
\end{remark}

\begin{remark}[What has not been measured at any sample size]
The transported cutoff bracket of Section~\ref{sec:chain} with its leakage
integral and the tracked decay rates; the per-stage sign condition as the
smallest eigenvalue of the increment compressed to the residual's subspace;
the floors of Section~\ref{sec:held-rough} in their Krylov form at a sample
size with a genuine late stage; and the
preconditioner term of the stage identity under Adam. Each is listed with
its decision criterion in the experimental plan accompanying this draft.
\end{remark}
